纳米硒（SeNPs）因具有良好的生物相容性和多重生物活性而受到广泛关注，但其毒性效应受粒径、浓度、表面修饰和材料化学状态等因素共同影响，传统实验方法难以快速完成系统评估。本文围绕纳米硒细胞毒性预测问题，构建了“基准数据集模型比较 + 专用数据集增强特征建模”的研究框架。首先，在来源于 Masarkar 等人的多类型纳米颗粒公开数据集 \cite{masarkar2024} 上，对随机森林、XGBoost、ABRF、LARF、TabTransformer 和蒸馏 TabTransformer 进行了系统比较；随后基于我们整理的 383 条硒纳米颗粒有效记录，进一步开展专用数据集特征统计、数据预处理分析与模型构建。

结果表明，在大样本异质表格数据上，TabTransformer 在回归任务中取得最高 $R^2$（0.768），而树模型在分类任务中表现更稳定，其中 XGBoost 的 AUC 为 0.945。在严格的硒纳米颗粒重建中，434 条记录按 24 个文献组进行 5 折外层、3 折内层嵌套交叉验证。无 IC50 的随机森林基础模型 OOF MAE 为 19.11\%、RMSE 为 24.12\%、$R^2=0.109$，XGBoost 分别为 19.49\%、23.97\% 和 0.120。化学增强和显式 log 剂量--log 时间交互均未产生稳定改善；IC50 辅助 XGBoost 的数值指标最好（MAE 18.35\%、RMSE 23.25\%、$R^2=0.172$），但其配对改进区间跨越 0，因此仅作为响应相关辅助参考。

固定预算的 General 到 SeNP 迁移对照同样未形成稳定收益：XGBoost 迁移的 MAE 点估计下降 0.67 个百分点，但按文献组 bootstrap 区间跨越 0，随机森林迁移的 MAE 反而上升。

上述结果表明，历史随机划分得到的高指标不能直接作为当前小样本异质归档数据的前瞻精度估计。现有可复现数据与评价流程可支持后续冻结模型外部验证，但在作出部署结论前仍需补齐修饰分类、分子量、配方标识及独立湿实验剂量--反应数据。
